\documentclass[11pt]{article}
\usepackage[utf8]{inputenc}
\usepackage[T1]{fontenc}
\usepackage{amsmath,amssymb,amsthm}
\usepackage[margin=2.5cm]{geometry}
\usepackage{hyperref}

\newtheorem{theorem}{Theorem}
\newtheorem{corollary}[theorem]{Corollary}
\newtheorem{definition}{Definition}

\title{BCCS: Formal Definition and NP-Hardness Proof\\[6pt]
\large Supplementary Material for ICDM 2026 Submission}
\author{Andre da Costa Silva}
\date{\today}

\begin{document}
\maketitle

%% ---------------------------------------------------------------
\section{Problem Definition}
%% ---------------------------------------------------------------

\begin{definition}[Budgeted Connected Case Segmentation --- BCCS]
\label{def:bccs}
Let $G=(V,E)$ be an undirected graph, $L\subseteq E$ a set of
\emph{target edges}, and $B\in\mathbb{Z}^+$ a \emph{budget}.
A \textbf{case assignment} is a partition
$\mathcal{P}=\{C_1,\dots,C_m\}$ of~$V$ into pairwise-disjoint,
non-empty subsets called \emph{cases}.

For each case~$C_i$ define the \textbf{induced edge set}
\[
  E(C_i)=\bigl\{(u,v)\in E : u\in C_i \text{ and } v\in C_i\bigr\}.
\]
The \textbf{coverage} of~$\mathcal{P}$ is
\[
  \operatorname{Cov}(\mathcal{P})
  =\bigl|\{(u,v)\in L:\exists\, C_i\in\mathcal{P}
            \text{ with } u,v\in C_i\}\bigr|.
\]
The \textbf{BCCS optimisation problem} is
\[
  \max_{\mathcal{P}}\;\operatorname{Cov}(\mathcal{P})
  \qquad\text{s.t.}\quad |E(C_i)|\le B
  \;\;\forall\, C_i\in\mathcal{P}.
\]
The \textbf{BCCS decision problem} asks: given $(G,L,B,\tau)$, does
there exist a partition~$\mathcal{P}$ with
$|E(C_i)|\le B$ for all~$i$ and
$\operatorname{Cov}(\mathcal{P})\ge\tau$?
\end{definition}

\paragraph{Interpretation in AML transaction monitoring.}
In anti-money laundering (AML) systems, $G$~represents the full
transaction graph, $L$~is the set of top-$k$ highest-risk edges
identified by a GNN scorer, and $B$~limits the number of edges an
analyst can review in a single investigation case.  BCCS asks how to
group transactions into bounded investigation cases to maximise the
number of flagged edges that are actually reviewable.

%% ---------------------------------------------------------------
\section{NP-Completeness}
%% ---------------------------------------------------------------

\begin{theorem}\label{thm:npc}
BCCS is NP-complete.
\end{theorem}

\begin{proof}

\subsection*{Membership in NP}

Given a partition $\mathcal{P}=\{C_1,\dots,C_m\}$, we can verify in
$O(|V|+|E|)$ time that (i)~$|E(C_i)|\le B$ for every~$i$
and (ii)~$\operatorname{Cov}(\mathcal{P})\ge\tau$, by building a
node-to-case lookup table and scanning all edges once.
Hence $\text{BCCS}\in\text{NP}$.

\subsection*{NP-Hardness (reduction from Maximum Independent Set)}

We reduce from the \textsc{Maximum Independent Set} problem (MIS),
known to be NP-hard~\cite{karp1972}.

\medskip\noindent
\textbf{MIS.}\;
Given a graph $H=(V_H,E_H)$ and an integer~$k$, does~$H$ contain an
independent set of size~$\ge k$?

\medskip\noindent
\textbf{Construction.}\;
Given a MIS instance $(H,k)$, build the BCCS instance
$(G,L,B,\tau)$ as follows.

\begin{enumerate}
\item \emph{Add a goal node.}
      Let~$g$ be a fresh node; set $V=V_H\cup\{g\}$.

\item \emph{Edges.}
      $E = E_H \;\cup\; \{(v,g) : v\in V_H\}$.

\item \emph{Target edges.}
      $L = \{(v,g) : v\in V_H\}$.
      Only the edges to~$g$ are targets.

\item \emph{Budget and threshold.}
      $B = \tau = k$.
\end{enumerate}

The construction runs in $O(|V_H|+|E_H|)$ time.

\medskip\noindent
\textbf{Intuition.}\;
Covering a target edge $(v,g)$ requires~$v$ and~$g$ in the same case.
Because~$g$ belongs to a single case, all ``covered'' nodes share that
case with~$g$.  The non-target edges from~$E_H$ between those nodes
consume budget, so only independent sets fit within budget~$k$.

\medskip\noindent
\textbf{Forward direction ($\Rightarrow$).}\;
Suppose~$H$ has an independent set $S\subseteq V_H$ with $|S|=k$.
Define the partition $C_1=S\cup\{g\}$; every remaining node forms a
singleton case.

\begin{itemize}
\item \emph{Budget of~$C_1$:}\;
      $E(C_1)=E_H(S)\cup\{(v,g):v\in S\}$.
      Since~$S$ is independent, $E_H(S)=\emptyset$, so
      $|E(C_1)|=|S|=k=B$.\;\checkmark

\item \emph{Singletons:}\; $0$~edges each.\;\checkmark

\item \emph{Coverage:}\;
      Each $v\in S$ is co-located with~$g$, giving
      $\operatorname{Cov}=k=\tau$.\;\checkmark
\end{itemize}

\medskip\noindent
\textbf{Backward direction ($\Leftarrow$).}\;
Suppose a partition~$\mathcal{P}$ achieves
$\operatorname{Cov}(\mathcal{P})\ge k$.
Let~$C_g$ be the case containing~$g$ and set
$S=C_g\cap V_H$.

A target edge $(v,g)$ is covered if and only if $v\in S$, so
$\operatorname{Cov}(\mathcal{P})=|S|\ge k$.

The induced edges of~$C_g$ are
\[
  E(C_g) = E_H(S)\;\cup\;\{(v,g):v\in S\},
\]
hence $|E(C_g)|=|E_H(S)|+|S|$.
The budget constraint gives
\[
  |E_H(S)|+|S|\;\le\; B = k.
\]
Since $|S|\ge k$, we obtain $|E_H(S)|\le k-|S|\le 0$.
Therefore $|E_H(S)|=0$ and $|S|=k$:
the set~$S$ is an independent set of size~$k$ in~$H$.\;\checkmark

\medskip
This completes the reduction.
\end{proof}

%% ---------------------------------------------------------------
\section{Consequences}
%% ---------------------------------------------------------------

\subsection{Inapproximability}

The reduction preserves the objective value exactly
($\operatorname{Cov}=|S|$).  Any $\alpha$-approximation for BCCS
therefore yields an $\alpha$-approximation for MIS.
Zuckerman~\cite{zuckerman2007} showed that MIS cannot be approximated
within $|V|^{1-\varepsilon}$ for any $\varepsilon>0$ unless P\,=\,NP.

\begin{corollary}\label{cor:inapprox}
Unless $\mathrm{P}=\mathrm{NP}$, BCCS cannot be approximated within
$|V|^{1-\varepsilon}$ for any $\varepsilon>0$.
\end{corollary}

This worst-case bound applies to adversarial instances.
Real-world AML graphs exhibit temporal locality and natural community
structure that make the problem far more tractable in practice.

\subsection{Connection to Maximum Coverage}

BCCS also relates to the \textsc{Maximum Coverage Problem} (MCP).
In settings where each case acts as a ``set'' of edges it covers,
the BCCS optimiser must decide which target edges to cover within
the budget, analogous to selecting sets in MCP.
Feige~\cite{feige1998} showed that MCP cannot be approximated beyond
$1-1/e\approx 0.632$ unless P\,=\,NP.

\subsection{Practical Implications}

Despite worst-case NP-hardness, our hierarchical WCC\,+\,Leiden
heuristic (BTCS) achieves \textbf{94--97\%} of the upper bound
(NoBudget WCC coverage) at $k=1\%$ on the AML100k and AML1M datasets.
This suggests that real-world AML transaction graphs possess structure
--- e.g.\ natural community boundaries that align with the budget ---
making the problem tractable in practice even though worst-case
instances are intractable.

%% ---------------------------------------------------------------
\begin{thebibliography}{9}

\bibitem{karp1972}
R.\,M.~Karp,
``Reducibility among combinatorial problems,''
in \emph{Complexity of Computer Computations},
pp.~85--103, Plenum Press, 1972.

\bibitem{feige1998}
U.~Feige,
``A threshold of $\ln n$ for approximating set cover,''
\emph{Journal of the ACM}, vol.~45, no.~4, pp.~634--652, 1998.

\bibitem{zuckerman2007}
D.~Zuckerman,
``Linear degree extractors and the inapproximability of max clique
and chromatic number,''
\emph{Theory of Computing}, vol.~3, no.~1, pp.~103--128, 2007.

\bibitem{traag2019}
V.\,A.~Traag, L.~Waltman, and N.\,J.~van Eck,
``From Louvain to Leiden: guaranteeing well-connected communities,''
\emph{Scientific Reports}, vol.~9, no.~1, p.~5233, 2019.

\end{thebibliography}

\end{document}
